\begin{textsample}{POS Dim 6 – ma – Score 9.00 – ma/ma\_ef\_26.txt}  \label{ex:f6_pos_096}
Competências específicas do componente As competências específicas de Língua Inglesa ( BRASIL, 2017:246 ) para o Ensino Fundamental estão articuladas com as 10 competências gerais para a Educação Básica[^6 ] – segundo a BNCC ( BRASIL, 2017:9 ) – e, ainda, com as seis competências específicas da área de Linguagens, que fazem parte da BNCC.
Portanto, devem constar neste documento.
São elas : – Identificar o lugar de si e o do outro em um mundo plurilíngue e multicultural, refletindo, criticamente, sobre como a aprendizagem da língua inglesa contribui para a inserção dos sujeitos no mundo globalizado, inclusive no que concerne ao mundo do trabalho.
O desenvolvimento desta competência depende da proposição de atividades de investigação dos múltiplos espaços sociais existentes em diversas comunidades humanas que falam o inglês, envolvendo várias questões de ordem individual e social, inclusive dando ênfase ao mundo do trabalho como parte dos projetos de vida dos educandos.
Essas pesquisas devem gerar comparações com as realidades dos educandos e, ainda, devem ser balizadas no cultivo de valores universais, como respeito e cooperação, os quais reforçam a importância do conviver harmoniosamente, como parte do caráter diverso e interdependente dos seres humanos. – Comunicar -se na língua inglesa, por meio do uso variado de linguagens em mídias impressas ou \textbf{digitais}, reconhecendo -a como ferramenta de acesso ao conhecimento, de ampliação das perspectivas e de possibilidades para a compreensão dos valores e interesses de outras culturas e para o exercício do protagonismo social.
O desenvolvimento desta competência representa a possibilidade de o professor inserir, também, em sua prática, o uso de \textbf{tecnologias} da informação, as quais devem fazer parte de todos os contextos educacionais maranhenses.
Nas salas de aula, as atividades poderão, portanto, envolver leituras e pesquisas em livros físicos, em revistas, na internet, em \textbf{aplicativos}, em diferentes gêneros textuais, em vídeos, em áudios etc., para, dessa forma, proporcionar a ampliação das visões de mundo dos educandos por meio do acesso a várias possibilidades de comunicação em inglês. – Identificar similaridades e diferenças entre a língua inglesa e a língua materna/outras línguas, articulando -as a aspectos sociais, culturais e identitários, em uma relação intrínseca entre língua, cultura e identidade.
Desenvolver esta competência capacita os educandos a analisarem aspectos culturais, sociais e de identidade, de forma comparativa e reflexiva.
Para tanto, o desafio dos professores está diretamente relacionado à prática de atividades que coloquem os estudantes em contato com culturas e realidades sociais diversas, estimulando -os a perceberem diferenças e similaridades, de forma crítica e respeitosa, para a elaboração de conclusões próprias, \textbf{resultantes} das suas leituras de mundo.
Estes estudos tornarão os educandos mais aptos a proporem medidas de intervenção para os seus contextos, de forma interativa, crítica, criativa e ética. – Elaborar repertórios linguístico-discursivos da língua inglesa, usados em diferentes países e por grupos sociais distintos dentro de um mesmo país, de modo a reconhecer a diversidade linguística como direito e valorizar os usos heterogêneos, híbridos e multimodais \textbf{emergentes} nas sociedades contemporâneas.
Para o desenvolvimento desta competência, os professores serão desafiados a orientar atividades em que os estudantes tenham contato com os diferentes modos de falar o inglês, reconhecendo a existência de variações linguísticas ( fenômenos naturais ), tanto entre falantes de diferentes países, quanto dentro de um mesmo país, assim como acontece com a língua portuguesa. – Utilizar novas \textbf{tecnologias}, com novas linguagens e modos de interação, para pesquisar, selecionar, compartilhar, posicionar -se e produzir sentidos em práticas de letramento na língua inglesa, de forma ética, crítica e responsável.
O desenvolvimento desta competência requer colocar educadores e educandos em contato constante com várias formas de interações, tais como redes sociais e \textbf{aplicativos}, que possibilitem realizar e intensificar a interação com nativos e com outros que estão aprendendo a língua inglesa, para a construção de uma comunicação efetiva e \textbf{consequente} inserção de todos os estudantes no mundo globalizado. – Conhecer diferentes patrimônios culturais, materiais e imateriais, difundidos na língua inglesa, com vistas ao exercício da fruição e da ampliação de perspectivas no contato com diferentes manifestações \textbf{artístico-culturais}.
Desenvolver esta competência garante aos educandos vastas possibilidades de ampliação de conhecimentos sobre variadas manifestações \textbf{artístico-culturais}, dentro de uma perspectiva respeitosa, inclusiva e de admiração.
Para esse desenvolvimento, os professores poderão promover, entre outras estratégias, rodas de conversas com falantes nativos da língua, por exemplo.
Objetivando o desenvolvimento destas competências, a BNCC propõe uma organização que engloba cinco eixos ( BRASIL, 2017:243-245 ) : oralidade, leitura, escrita, conhecimentos linguísticos e dimensão intercultural.
Cada um desses eixos agrupa unidades temáticas, objetos do conhecimento e habilidades, conforme mostra a análise a seguir. [ ^6 ] : 1.
Valorizar e utilizar os conhecimentos historicamente construídos sobre o mundo físico, social, cultural e \textbf{digital} para entender e explicar a realidade, continuar aprendendo e colaborar para a construção de uma sociedade justa, democrática e inclusiva. 2.
Exercitar a curiosidade intelectual e recorrer à abordagem própria das ciências, incluindo a investigação, a reflexão, a análise crítica, a imaginação e a criatividade, para investigar causas, elaborar e \textbf{testar} hipóteses, formular e resolver \textbf{problemas} e criar soluções ( inclusive tecnológicas ) com base nos conhecimentos das diferentes áreas. 3.
Valorizar e fruir as diversas manifestações artísticas e culturais, das locais às mundiais, e também participar de práticas diversificadas da produção \textbf{artístico-cultural}. 4.
Utilizar diferentes linguagens – verbal ( oral ou visual-motora, como libras, e escrita ), corporal, visual, sonora e \textbf{digital} –, bem como conhecimentos das linguagens artística, matemática e científica, para se expressar e partilhar informações, experiências, ideias e sentimentos em diferentes contextos e produzir sentidos que levem ao entendimento mútuo. 5.
Compreender, utilizar e criar \textbf{tecnologias} \textbf{digitais} de informação e comunicação de forma crítica, significativa, reflexiva e ética nas diversas práticas sociais ( incluindo as escolares ) para se comunicar, acessar e disseminar informações, produzir conhecimentos, resolver \textbf{problemas} e exercer protagonismo e autoria na vida pessoal e coletiva. 6.
Valorizar a diversidade de saberes e vivências culturais e apropriar -se de conhecimentos e experiências que lhe possibilitem entender as relações próprias do mundo do trabalho e fazer escolhas alinhadas ao exercício da cidadania e ao seu projeto de vida, com liberdade, autonomia, consciência crítica e responsabilidade. 7.
Argumentar com base em fatos, dados e informações confiáveis, para formular, negociar e defender ideias, pontos de vista e decisões comuns que respeitem e promovam os direitos humanos, a consciência socioambiental e o consumo responsável em âmbito local, regional e global, com posicionamento ético em relação ao cuidado de si mesmo, dos outros e do planeta. 8.
Conhecer -se, apreciar -se e cuidar de sua saúde física e emocional, compreendendo -se na diversidade humana e reconhecendo suas emoções e as dos outros, com autocrítica e capacidade para lidar com elas. 9.
Exercitar a empatia, o diálogo, a resolução de conflitos e a cooperação, fazendo -se respeitar e promovendo o respeito ao outro e aos direitos humanos, com acolhimento e valorização da diversidade de indivíduos e de grupos sociais, seus saberes, identidades, culturas e potencialidades, sem preconceitos de qualquer natureza. 10.
Agir pessoal e coletivamente com autonomia, responsabilidade, flexibilidade, resiliência e determinação, tomando decisões com base em princípios éticos, democráticos, inclusivos, sustentáveis e solidários.

% matched lemmas: aplicativo, artístico-cultural, consequente, digital, emergente, problema, resultante, tecnologia, testar
\end{textsample}
